\section{Conductor three in the actual character-supply interface}
\label{sec:conductor-three-audit}

\begin{proposition}[AC-G34: conductor three is allowed by the inspected supply chain]
\label{AC-G34}
In the supplied character-supply interface, the primitive character
$\psi_{-3}$ is not excluded by the separate predicate
\texttt{CoprimeFifteen}.  More precisely, the declarations and
implementations listed below accept an arbitrary primitive character
$\psi$ of conductor $M>0$; the relevant coprimality condition is
$\gcd(M,qd)=1$.  Specialization to $M=3$, with the large auxiliary
primes $q,d$ in the supplied windows, satisfies this condition.
The conclusion remains conditional on the actual supplied
$\operatorname{VWinAt}(\psi,401)$ input.  It is the pure Tate
Selmer bound in the source, not a cubic Artin twist of that bound.

The source also contains a precise reserve for a \emph{fixed}
frame multiplier: it passes from its cutoff $\lfloor x'^{401}\rfloor$
to $\lfloor x^{400}\rfloor$ under a constant rescaling
$x=kx'$.  This reserve does not establish the independent-window
amalgamation considered in AC-G31.
\end{proposition}
\begin{proof}
\proofstep{1. Exact declarations and the conductor.}
All fetched files in this audit are pinned to repository commit
\texttt{7cbf7a0c32fca1badf05931879a37c3f8a2b03b4}; their complete
contents and hashes are retained in
\path{evidence/source_interface/manifest.json}.
The exact outer declaration is
\texttt{UpperTheory.Top.supplyChi} in
\texttt{CharacterBoxes/SupplyChi.lean}:
\[
 \forall\psi:\mathrm{PChar},\qquad
 \operatorname{VWinAt}(\psi,401)
       \Longrightarrow
 \operatorname{HostExW}(3/2,(\operatorname{GChi}\psi).\mathrm{req}).
                                                        \tag{G34.1}
\]
The implementation calls
\texttt{UpperTheory.T3.supplyOfAdmOf}, in the file of that name,
and the admission theorem in
\texttt{CoprimeRanks/AdmChi.lean}.
The character-field construction is the theorem
\texttt{UpperTheory.Top.fieldChi} in
\texttt{CharacterBoxes/FieldChi.lean}.
Its hypotheses on the auxiliary primes include
\[
 q\ne3,\quad \gcd(M,qd)=1,\quad 4M\le b,
 \qquad
 \left(\frac{-d}{\ell}\right)=1
 \quad(\ell\le b,\ \ell\text{ odd prime}),             \tag{G34.2}
\]
and its conductor conclusion is
\[
 M\mid Q_S,\qquad
 Q_S\mid\operatorname{lcm}(2^{e+2},M),\qquad
 Q_S\le4M\,\operatorname{ord}(\phi).                   \tag{G34.3}
\]
The first divisibility is the first clause of
\texttt{KerChi}, defined in \texttt{CharacterBoxes/TwistForm.lean}.
For $M=3$, (G34.3) therefore permits, and indeed requires,
$3\mid Q_S$.  The condition $q\ne3$ in (G34.2) is a condition
on the auxiliary prime, not on the character conductor.
In \texttt{supplyOfAdmOf}, it is proved from the large-prime window;
$\gcd(M,qd)=1$ is supplied by \texttt{PairGood.coprime}.
Thus there is no conductor-three exclusion in this step.

The theorem \texttt{UpperTheory.T3.hostOfFieldOf} in
\texttt{ShortNorms/hostOfFieldOf.lean} accepts $M>0$ with
$Q_S\mid\operatorname{lcm}(2^{e+2},M)$ and does not assume
$\gcd(Q_S,15)=1$.  In its helper \texttt{hof\_exists\_host},
the congruences modulo fifteen concern the chosen level
primes $p_{\mathrm{de}},p'_{\mathrm{de}}$:
\[
 p_{\mathrm{de}}\equiv p'_{\mathrm{de}}\equiv1\pmod{15}.
\]
They are obtained together with congruence one modulo
$Q_Sqd$, using a common multiple.  Both congruences are
compatible even when three divides $Q_S$.
The separately named \texttt{CoprimeFifteen}, whose definition
is \texttt{Nat.Coprime H.QS 15} in
\texttt{ConjEMemHostFIff.lean}, is neither a field of
\texttt{Host} nor a hypothesis of the displayed chain.
The full implementations of \texttt{FieldChi},
\texttt{WinPairs}, \texttt{supplyOfAdmOf},
\texttt{hostOfFieldOf}, \texttt{PairRows}, \texttt{PairHyp},
\texttt{IsChiField} and \texttt{AdmChi} were read for this audit.
This proves the claimed scope of the conductor check.

\proofstep{2. The fixed rescaling and its exact cutoff.}
Put $a=2+8\varepsilon$ and let the fixed multiplier be
$c\ge1$.  Set $k=c^{1/a}$, so $k\ge1$ and $k^a=c$.
The dyadic page at $x'$ has
$x'^a\le2^e\le2x'^a$; therefore, at $x=kx'$,
\[
                  x^a\le c2^e\le2x^a.                \tag{G34.4}
\]
If $x'\ge k^{400}$, then
\[
 x^{400}=k^{400}x'^{400}\le x'^{401},\qquad
 \sqrt{\exp(x'^\varepsilon)}
        \le\sqrt{\exp(x^\varepsilon)}.                \tag{G34.5}
\]
These inequalities give the required monotonicity of the
Selmer cutoff and budget.  If the old window is $[W,W^2]$
and $W\ge k^3$, then
$(kW)^{3/2}\le W^2$; every $x$ in
$[kW,(kW)^{3/2}]$ has $x'=x/k$ in the old window.
This is exactly the argument of
\texttt{UpperTheory.Top.winPairs}.
Its target \texttt{WinDatum.bound}, defined in
\texttt{TwistForm.lean}, is
\[
 \operatorname{VVle}
   \bigl(E,\lfloor x^{400}\rfloor,
                 \sqrt{\exp(x^\varepsilon)}\bigr),    \tag{G34.6}
\]
for every field satisfying the literal
\texttt{IsChiField} condition.
The subsequent source rank theorem bounds its chosen $R.t$
by that cutoff.  Thus one must distinguish the outer
$401$-cutoff hypothesis from the $400$-cutoff reserve used
in constructing a ranked host.

No exponent $\varepsilon$ is changed in this argument, and
there is only one character window.  In particular it does
not remove the two-scale $q,d$ obstruction in AC-G31.
Finally, the conclusion of \texttt{supplyOfAdmOf} explicitly
uses \texttt{VVle H.E R.t (Real.sqrt s.b)}; the coefficient
there is the pure Tate one in \texttt{LData}, as in AC-G29.
Replacing it by a nontrivial cubic coefficient is a different
statement.  This proves all assertions.
\end{proof}

This is a check of the supplied source statements and their
inspected mathematical implementation, not a claim that a
new Lean build was executed.  The audit closes the proposed
\texttt{CoprimeFifteen} exclusion at this interface; it neither
derives \texttt{VWinAt} from theorem A nor changes the
coefficient in \texttt{NData.sha\_eq}.
